% META: {"id": "1SPE-PROBCOND-FR-R5", "chapitre": "1SPE-PROBA-COND", "type_objet": "remediation", "capacites_codes": ["R5"], "status": "generated"}

%---------------------------------------------------------------
% FICHE DE REMISE A NIVEAU --- R5 : Tableaux de donnees
%---------------------------------------------------------------

\subsection*{Rappel --- Tableaux de donnees}

Un \textbf{tableau a double entree} croise deux criteres. La \textbf{frequence} d'une categorie est le rapport de son effectif au total.

\begin{exercice}{1SPE-PROBCOND-FR-R5-EX1}{1}{5}
\begin{center}
\begin{tabular}{|c|c|c|c|}
\hline
 & Bus & Velo & Total \\
\hline
Garcons & $30$ & $20$ & $50$ \\
\hline
Filles & $40$ & $10$ & $50$ \\
\hline
Total & $70$ & $30$ & $100$ \\
\hline
\end{tabular}
\end{center}

\begin{enumerate}
  \item Quelle est la frequence des cyclistes ?
  \item Quelle est la frequence des filles parmi les cyclistes ?
\end{enumerate}
\end{exercice}

% BEGIN-VERIFY
% from sympy import Rational
% assert Rational(30, 100) == Rational(3, 10)
% assert Rational(10, 30) == Rational(1, 3)
% END-VERIFY

\begin{corrige}{1SPE-PROBCOND-FR-R5-EX1}
\textbf{1.} Frequence des cyclistes : $\frac{30}{100} = \frac{3}{10} = 30\,\%$.

\textbf{2.} Parmi les cyclistes : $\frac{10}{30} = \frac{1}{3} \approx 33{,}3\,\%$.
\end{corrige}
